\chapter{Conclusion}\label{ch:conclusion}

This dissertation began with a narrow, concrete complaint: that porting a parallel
algorithm across CPU threads, a message-passing cluster, and an embedded
microcontroller means rewriting it, because the decomposition and the IO --- not the
arithmetic --- are written by hand and tangled into the source, separately for each
target. Compute scheduling has long had first-class language support; IO scheduling
has not. The work set out to give it that support and to make the resulting
separation claim something stronger than a design aspiration.

The contribution is best stated soberly. The system rests on a strict separation
between an algorithm sublanguage, which says what to compute and contains no worker,
buffer, transport, or synchronisation, and a schedule sublanguage, in which
decomposition, IO semantics, and target are first-class directives the algorithm
never sees. The two are unified by a single Petri-net intermediate representation
that carries transfer synthesis, synchronisation, and a compile-time soundness gate
together, so that boundedness and deadlock-freedom are reported as compile errors
rather than discovered as runtime hangs. And the genericity claim is made
mechanically falsifiable by a cross-backend differential test: the same algorithm,
under every schedule and capable backend, must produce byte-identical output, and
any leak across the algorithm/schedule or middle-end/presentation boundary that
makes two backends diverge turns that test red. The implementation drives
twenty-seven worked algorithms through ten
backends across three tiers --- seven commodity-CPU backends, two message-passing
cluster backends, and one embedded backend with multi-microcontroller co-simulation.

What was demonstrated is correspondingly bounded. The seven CPU backends agree
byte-for-byte across the exercised schedule-by-backend matrix; the soundness gate
accepts every sound schedule in the corpus and is shown, by adversarial tests, able
to reject unsound nets; every check that underwrites these claims has a negative
test proving it can fail, so a green result is informative rather than vacuous; and
the cluster and embedded tiers are confirmed to value-correctness on a local
message-passing launcher and to byte-exact co-simulation respectively. None of this
is a proof that every program the grammar admits is byte-identical everywhere; it is
the result of a falsification rig that has had every opportunity over a broad and
diverse corpus to refute the central claim and has not, with its documented edges
marking exactly where the claim has not yet been tested.

The work is equally clear about its price. The whole architecture is bought with a
single trade --- expressiveness for a statically determined, deterministic firing
order --- and that trade defines a real boundary. The algorithm class is affine,
static, single-assignment, and integer-valued, with floating-point summation
excluded by the determinism it would break; there is no data-dependent control flow,
no recursion, and no dynamic scheduling; the schedule is checked, not optimised, and
carries no cost model; the soundness guarantee is scoped to coordination within a
restricted net class, not to value-correctness or to arbitrary nets; and the cheap
CPU tier carries the strongest form of the formal claim while the more exotic tiers
are confirmed more narrowly. These are consequences of the central commitment, not
oversights, and the dissertation has tried throughout to weigh them with the same
seriousness as the advantages.

The bet the work makes is that for a particular and not uncommon audience --- people
moving regular, deterministic algorithms across heterogeneous parallel and embedded
targets, who value portability and provable coordination over peak performance and
can live within the affine, static class --- a system that makes IO a first-class
schedule concern, decides soundness at compile time, and renders its genericity
claim falsifiable is worth more than a more expressive system whose portability is
asserted rather than tested. Whether that bet generalises beyond the demonstrated
class is the substance of the future work --- including order-independent
floating-point reduction, bounded data-dependent iteration, precise distributed
communication, further hardware families, prescriptive scheduling, and the deferred
runtime performance study --- each of which would widen the audience or sharpen the
evidence. The discipline that
shaped the original design --- lift a limit only while keeping the soundness gate
sound and the differential test meaningful --- remains the constraint under which
that widening must happen.

The smallest honest summary is this. The expensive part of portable parallel
software is not the arithmetic; it is the decomposition and the IO. This work
factored those out into a separate, swappable schedule, unified them in one analysable
representation, and then --- rather than asking to be believed --- built the
instrument that would expose the separation if it leaked, and ran it on every build.
Across the corpus, it did not leak. That is a modest claim, precisely scoped, and it
is the one the dissertation stands behind.
